% Image credits for the photographs and AI illustrations of Book 4 (University, Year 2).
% Machine-readable license records live in images/book4/CREDITS.md; keep the
% two files in sync. The Book 4 writer owns this file: add one \item per
% photograph (inside the itemize, by chapter) as photographs land.
\chapter*{Bronvermelding van afbeeldingen}

De tekeningen, kaarten en grafieken in dit boek zijn origineel; kaarten en
grafieken die op gepubliceerde gegevens zijn gebaseerd, vermelden hun bron in
het bijschrift of in deze lijst. Foto's worden gebruikt onder hun respectieve
vrije licenties, met dank aan hun makers; volledige bronlinks en
licentie-URL's voor elke foto staan in \texttt{images/book4/CREDITS.md} in de
repository van het boek.

\bigskip
Naast de foto's en de originele tekeningen zijn sommige figuren door AI
gegenereerde illustraties, gemaakt met OpenAI-beeldgeneratie op basis van
prompts die door de redacteuren van het boek zijn geschreven, en vóór opname
gecontroleerd op wetenschappelijke juistheid. Ze tonen uitsluitend generieke
scènes, nooit een echte plaats, een echt persoon of een echt exemplaar. De
volledige prompt voor elke gegenereerde afbeelding staat in
\texttt{images/book4/ai/PROMPTS.md} in de repository.

\bigskip
\noindent Foto's en gegevens, per hoofdstuk:
\begin{itemize}\small
  \item Hfst.~1 (het interstellaire medium en de vorming van sterren): de protoplanetaire schijf van HL Tauri, ALMA (ESO/NAOJ/NRAO), CC BY 4.0; de Pilaren van de Schepping in de Adelaarnevel, NASA, ESA en het Hubble Heritage Team (STScI/AURA), publiek domein.
  \item Hfst.~2 (de vorming van planeten en meteorieten): de koolstofhoudende chondriet Allende en het Widmanstättenpatroon van de ijzermeteoriet Gibeon, James St.~John, CC BY 2.0.
  \item Hfst.~3 (de elementen in de aarde): een xenoliet van lherzoliet uit de mantel in basalt, Kilbourne Hole, James St.~John, CC BY 2.0.
  \item Hfst.~5 (magma's en hun intrusie): een gang bij Shiprock, Dexter Perkins, CC0 (Wikimedia Commons).
  \item Hfst.~6 (vulkaanuitbarstingen): een lavafontein op de kegel Pu`u~`\=O`\=o van de Kīlauea, J.\,D. Griggs, USGS, publiek domein; de eruptiekolom van de Mount St.~Helens, 18 mei 1980, Austin Post, USGS, publiek domein; de eruptiekolom van de Pinatubo, 12 juni 1991, Dave Harlow, USGS, publiek domein.
  \item Hfst.~7 (de mantel smelten): een handstuk duniet, Mike Davis (GeoDIL), CC0.
  \item Hfst.~8 (magmatische differentiatie): chromitietlagen van het Bushveldcomplex, Kevin Walsh, CC BY 2.0; gelaagde gabbro's van de Skaergaard-intrusie, C. Tegner, CC BY 4.0.
  \item Hfst.~10 (metamorfe reacties en facies): lawsonietblauwschist, Ward Creek, Californië, en eclogiet, Almenningen, Noorwegen, James St.~John, CC BY 2.0.
  \item Hfst.~11 (thermobarometrie en P--T--t-paden): granaatschist, James St.~John, CC BY 2.0; een coësietkorrel in eclogiet, gekruiste polarisatoren, J.\,R. Smyth, CC BY-SA 1.0.
  \item Hfst.~12 (smelten van de korst: migmatieten en granieten): een ontsluiting van migmatiet in Nashua, New Hampshire, Kevin M. Gill, CC BY 4.0; een leucogranietader, Anne Burgess (Geograph), CC BY-SA 2.0.
  \item Hfst.~13 (anatomie van een gebergtegordel: de Alpen): kussenlava's van de ofioliet van de Chenaillet, Lionel Duchoiselle, CC BY-SA 3.0; de Matterhorn vanuit Zermatt, Ximonic (Simo Räsänen), CC BY-SA 4.0.
  \item Hfst.~14 (botsingsgordels en Andesgordels vergeleken): het Tibetaans Plateau tussen Nepal en Shigatse, Jan Kranendonk, CC BY-SA 4.0; de Boliviaanse Altiplano, Luca Galuzzi, CC BY-SA 2.5.
  \item Hfst.~15 (oude massieven en de Wilsoncyclus): Siccar Point, zuidoostelijke klif, dave souza, CC BY-SA 4.0.
  \item Hfst.~16 (sedimentaire bekkens vergeleken): een halfafzinkbaar booreiland in de Noordzee, Erik Christensen, CC BY-SA 3.0.
  \item Hfst.~17 (oceaaneilanden): het atol Kure vanuit het Internationaal Ruimtestation, NASA, publiek domein.
  \item Hfst.~18 (de precambrische Aarde): spinifexkomatiiet bij Pyke Hill, Ontario, James St.~John, CC BY 2.0; kapcarbonaat van de Nuccaleena-formatie, Zuid-Australië, James St.~John, CC BY 2.0.
  \item Hfst.~19 (de fanerozoïsche Aarde): krijtpilaren van Monument Rocks, Kansas, James St.~John, CC BY 2.0; het niveau van Bonarelli, Picenus, CC BY-SA 4.0.
  \item Hfst.~20 (fossielen en paleomilieus): \emph{Waptia} en \emph{Marrella} uit de Burgess Shale, James St.~John, CC BY 2.0; het Berlijnse exemplaar van \emph{Archaeopteryx}, H.~Raab, CC BY-SA 3.0; \emph{Zoophycos} bij Sagres, Amada44, CC BY 4.0.
  \item Hfst.~21 (thermodynamica van de atmosfeer): het aambeeld van een supercelonweer, Stephen Corfidi, NOAA, publiek domein; radiosondepeilingen van het Integrated Global Radiosonde Archive, I. Durre en anderen, NOAA NCEI, publiek domein.
  \item Hfst.~22 (weersystemen): een lagedrukgebied boven IJsland, NASA/GSFC, MODIS Rapid Response Team, Jacques Descloitres, publiek domein; het oog van orkaan Isabel vanuit het Internationaal Ruimtestation, Ed Lu, NASA, publiek domein.
  \item Hfst.~23 (klimaatvariabiliteit): indices van de Noord-Atlantische Oscillatie en de Oceanic Niño Index, NOAA Climate Prediction Center, publiek domein.
  \item Hfst.~24 (klimaatarchieven en proxy's): jaarlagen van de GISP2-ijskern op 1855~m, Anthony Gow, U.S. Army Corps of Engineers CRREL (NOAA Paleoclimatology), publiek domein; doorsnede van een speleotheem van travertijn, Diamond Caverns, Kentucky, James St.~John, CC BY 2.0; de bentische LR04-stapel van L.\,E. Lisiecki en M.\,E. Raymo (2005), de temperatuurreeks van EPICA Dome C van J. Jouzel en anderen (2007) en de Antarctische \ce{CO2}-samenstelling van B. Bereiter en anderen (2015), NOAA NCEI Paleoclimatology, publiek domein.
  \item Hfst.~25 (orbitale forcering en de ijstijden): baanelementen en instraling van Berger en Loutre (1991), NOAA NCEI Paleoclimatology, publiek domein.
  \item Hfst.~26 (het klimaat over de geologische tijd): de bentische isotopenreeks van het Cenozoïcum CENOGRID, T. Westerhold en anderen (2020), PANGAEA, CC BY 4.0; de fanerozoïsche \ce{CO2}-compilatie van G.\,L. Foster, D.\,L. Royer en D.\,J. Lunt (2017), \emph{Nature Communications}, CC BY 4.0.
  \item Hfst.~27 (antropogene klimaatverandering): IPCC Sixth Assessment Report, Working Group~I, uitgebreide gegevens van Annex~III, CC BY 4.0.
  \item Hfst.~28 (geologische gevaren en risico): puinstroomafzettingen bij Caraballeda, Venezuela, Lawson Smith, U.S. Army Corps of Engineers, publiek domein; een moskee die na de tsunami van 2004 in Atjeh bleef staan, USGS, publiek domein.
  \item Hfst.~29 (ertsafzettingen): de Bingham Canyon Mine vanuit het Internationaal Ruimtestation, NASA (bemanning van Expedition~15), publiek domein; de ijzeren hoed van de Flambeau Mine, Wisconsin, James St.~John, CC BY 2.0; gehalte--tonnagegegevens van D.\,A. Singer, V.\,I. Berger en B.\,C. Moring (2008), USGS, publiek domein.
  \item Hfst.~30 (energiegeologie): een geothermische centrale bij Geyserville, Californië, U.S. Department of Defense, publiek domein; dichtheden van koolstofdioxide uit het NIST Chemistry WebBook (Standard Reference Database 69).
  \item Hfst.~31 (veldwerk- en laboratoriummethoden II): poikilitische augiet die plagioklaas omsluit, Strekeisen, CC BY-SA 4.0; slijpplaatje van een granaat--glimmerschist, Jackdann88, CC BY 3.0; elektronenmicrosonde bij de Geological Survey of Finland, Kallerna, CC BY-SA 4.0; silicaresultaten voor referentiematerialen van gesteenten, J.\,E. Taggart en D.\,F. Siems (2002), USGS, publiek domein.
\end{itemize}
